% ===============================================================
%  PART IV -- Our Contribution (Chapters 15--16)
%  Method and analysis only. Results are in Part VI; the simulator
%  is described in Part V. Proofs are in Appendix D (rule 7).
% ===============================================================

\newtheorem{proposition}{Proposition}[chapter]

\part{Our Contribution}
\label{part:contribution}

\begin{center}
\adjustbox{max width=\linewidth}{%
\begin{forest}
  part map,
  [Part IV
    [Ch.~15: A Controlled Method
      [One changed variable] [Everything else held fixed] [Paired runs] [Correctness first]]
    [Ch.~16: The Dirty Sharing Ratio
      [One counter under MESI] [Dirty epochs] [An exact saving identity] [From bytes to time] [Hypotheses H1--H4]]
  ]
\end{forest}}
\end{center}

\begin{keybox}[What Part IV does]
Part~III left two open questions (\cref{sec:cmp-gap}): none of the studies we
found isolates the effect of the \stO\ state, and none offers a workload measure,
taken under MESI, that predicts the gain from moving to MOESI. Chapter~15 answers
the first with a controlled comparison in which \stO\ is the only difference
between the two protocols. Chapter~16 answers the second with the Dirty Sharing
Ratio and derives an exact identity linking it to the traffic MOESI saves. Four
hypotheses, one per research question, close the Part. Everything here is derived
from the counting rules of \cref{tab:counting-rules}; no result depends on a
simulation.
\end{keybox}

% ===============================================================
%  Chapter 15 -- A Controlled Method for Isolating the Owned State
% ===============================================================
\chapter{A Controlled Method for Isolating the Owned State}
\label{ch:method}

\begin{figure}[htbp]
\centering
\adjustbox{max width=\linewidth}{%
\begin{forest}
  chapter map,
  [Chapter 15
    [Why comparisons mislead [Whole protocols differ] [Confounders]]
    [The one changed variable [Three rules differ] [Diff of the FSMs]]
    [Held fixed {\&} paired [Controlled variables] [Same trace twice]]
    [Same cache contents [Proposition 15.1] [Worked example]]
    [Correctness first [Invariants] [Version numbers]]
  ]
\end{forest}}
\caption[]{Structure of Chapter~15: why existing
comparisons cannot isolate the \stO\ state, the single change the method makes,
what it holds fixed, and how the runs are checked.}
\label{fig:ch15-map}
\end{figure}

% ---------------------------------------------------------------
\section{Why Existing Comparisons Cannot Isolate the Owned State}
\label{sec:method-why}
\sectionmap{[The difficulty [A protocol is many decisions] [Implementations differ in all of them]
  [A measured gap has many causes] [We need one variable]]}
\index{MESI}\index{MOESI}\index{Owned state}\index{coherence axioms}%

A coherence protocol is made of many decisions: which states exist, who supplies
data on a miss, when memory is updated, whether a clean block may be forwarded
between caches, how races are resolved and how each message is timed. Two
published implementations called ``MESI'' and ``MOESI'' differ in many of these at
once. Chapter~14 showed this for the evaluations we reviewed
(\cref{tab:evaluations}): the early study compared whole protocols, and the
simulator studies used implementations that differ in several states and
transitions, in one case without the \stM$\to$\stO\ transition at all. When such a
study reports that MOESI is faster, the gap could come from \stO, from a better
forwarding policy, from a different cache set-up, or from all of these.
\Cref{fig:method-idea} contrasts that kind of comparison with ours. Our method
removes each of these sources of difference: MOESI is MESI plus \stO\ with every
other rule identical (\cref{sec:method-delta}), both runs share one cache model, bus
and set of message sizes and replay the same trace (\cref{sec:method-fixed}),
differences are taken per trace before averaging over seeds, and both coherence
axioms are checked on every access so that a bug cannot pass for a speed-up
(\cref{sec:method-correct}).

\begin{figure}[htbp]
\centering
\begin{bigfig}[0.95]
\begin{tikzpicture}[font=\sffamily\scriptsize,
    bx/.style={draw, thick, rounded corners=2pt, minimum width=2.15cm, minimum height=0.55cm, align=center},
    ar/.style={-{Stealth}, thick}]
  \node[font=\sffamily\bfseries\small] at (2.6,4.65) {(a) Whole-protocol comparison};
  \foreach \y/\t in {3.8/States, 3.1/Transitions, 2.4/Forwarding, 1.7/Cache config., 1.0/Timing model} {
    \node[bx, fill=mesiBlue!12, draw=mesiBlue, minimum width=2.25cm] at (1.15,\y) {\t\ A};
    \node[bx, fill=moesiOrange!14, draw=moesiOrange, minimum width=2.25cm] at (4.05,\y) {\t\ B};
    \node[text=alertRed, font=\sffamily\bfseries] at (2.6,\y) {$\neq$};
  }
  \node[bx, fill=alertRed!10, draw=alertRed, minimum width=5.1cm] (ra) at (2.6,0.05) {gap in result: cause unknown};
  \draw[ar, alertRed] (1.15,0.72) -- (1.15,0.33); \draw[ar, alertRed] (4.05,0.72) -- (4.05,0.33);
  \begin{scope}[xshift=6.9cm]
  \node[font=\sffamily\bfseries\small] at (2.75,4.65) {(b) Controlled comparison (this report)};
  \node[bx, fill=coreGrey!10, draw=coreGrey, minimum width=5.5cm] at (2.75,3.8) {one trace, one interleaving};
  \node[bx, fill=coreGrey!10, draw=coreGrey, minimum width=5.5cm] at (2.75,3.1) {one cache model (size, LRU, 64-B blocks)};
  \node[bx, fill=coreGrey!10, draw=coreGrey, minimum width=5.5cm] at (2.75,2.4) {one bus, one set of message sizes};
  \node[bx, fill=mesiBlue!12, draw=mesiBlue, minimum width=2.3cm] (tm) at (1.15,1.5) {MESI rules for\\remote read of \stM};
  \node[bx, fill=moesiOrange!14, draw=moesiOrange, minimum width=2.3cm] (to) at (4.35,1.5) {MOESI rules for\\remote read of \stM};
  \draw[{Stealth}-{Stealth}, thick, memGreen] (tm.east) -- node[above, font=\sffamily\bfseries\tiny, text=memGreen]{swap} (to.west);
  \node[bx, fill=memGreen!12, draw=memGreen, minimum width=5.5cm] (rb) at (2.75,0.05) {gap in result: caused by \stO\ alone};
  \draw[ar, memGreen] (tm) -- (tm|-rb.north); \draw[ar, memGreen] (to) -- (to|-rb.north);
  \end{scope}
\end{tikzpicture}
\end{bigfig}
\caption[Whole-protocol versus controlled comparison]{(a) Comparing two separately
built protocols changes many things at once. (b) Our method runs one simulator
twice on the same input and changes only the rules that involve \stO.}
\label{fig:method-idea}
\end{figure}

% ---------------------------------------------------------------
\section{The One Changed Variable}
\label{sec:method-delta}
\sectionmap{[The delta [MOESI as MESI plus O] [Three rules change]
  [All other rules identical] [Written as a diff]]}
\index{MESI}\index{MOESI}\index{Owned state}\index{owner}%

MOESI is a class of protocols whose members may choose different allowed actions
(\cref{sec:own-moesi}), so a comparison has to fix one member. We take the member
closest to MESI: it adds \stO\ and changes only the three rules that \stO\ forces,
and is otherwise the MESI of \cref{fig:mesi}.
\begin{enumerate}[itemsep=1pt]
  \item Remote read of a block in \stM. MESI: the holder supplies the block,
    writes it back and moves to \stS. MOESI: the holder supplies the block without
    writing it back and moves to \stO.
  \item Remote read of a block in \stO\ (MOESI only): the owner supplies the block
    and stays in \stO; memory is not involved.
  \item Eviction of a block in \stO\ (MOESI only): the owner writes the block back,
    as an eviction from \stM\ does. \stS\ copies stay.
\end{enumerate}
A write by any cache to a block in \stO\ or \stS\ is an upgrade in both protocols:
one control message invalidates all other copies and the writer moves to \stM. If
the writer did not hold the block, the owner (in \stM\ or \stO) supplies it and is
invalidated, and memory is not written because the writer becomes the new dirty
owner. \Cref{fig:moesi-diff} draws the difference as a diff of the two state
machines and \cref{tab:transition-diff} lists every rule. Other members of the
class may also forward clean blocks or hand ownership to the latest reader, but
each of those is a second change. Leaving them out answers the narrow question of
what \stO\ itself does; such features could be added later and studied one at a
time with the same method.

\begin{figure}[htbp]
\centering
\begin{bigfig}[1.0]
\begin{tikzpicture}[fsm]
  \node[font=\sffamily\bfseries\small] at (2.5,4.5) {(a) MESI: remote read of \stM};
  \node[stB] (M) at (0,3.0) {\stM};
  \node[stB] (S) at (5.0,3.0) {\stS};
  \node[stB] (I) at (2.5,0.2) {\stI};
  \draw[key] (M) -- node[lb, above, text=alertRed]{readmiss / send data\\+ \textbf{write back}} (S);
  \draw[bus] (M) -- node[lb, left=3pt]{writemiss /\\send data} (I);
  \draw[bus] (S) -- node[lb, right=3pt]{write, writemiss /\\---} (I);
  \node[lb, text=coreGrey] at (2.5,1.95) {memory now up to date};
  \begin{scope}[xshift=7.6cm]
  \node[font=\sffamily\bfseries\small] at (2.5,4.5) {(b) MOESI: the same event};
  \node[stO] (M2) at (0,3.0) {\stM};
  \node[stO] (O2) at (5.0,3.0) {\stO};
  \node[stO] (I2) at (0,0.2) {\stI};
  \node[stO] (S2) at (5.0,0.2) {\stS};
  \draw[good] (M2) -- node[lb, above, text=memGreen]{readmiss / send data\\(no write back)} (O2);
  \draw[good] (O2) to[out=30, in=-30, looseness=7] node[lb, right]{readmiss /\\send data} (O2);
  \draw[bus] (M2) -- node[lb, left=3pt]{writemiss /\\send data} (I2);
  \draw[bus] (O2) to[bend right=12] node[lb, sloped, above, pos=0.42]{write, writemiss / send data if needed} (I2);
  \draw[tr] (O2) to[bend left=12] node[lb, sloped, below, pos=0.58]{evict / write back} (I2);
  \draw[bus] (S2) -- node[lb, below]{write, writemiss / ---} (I2);
  \end{scope}
\end{tikzpicture}
\end{bigfig}
\caption[MESI and MOESI as a diff]{The only transitions that differ between the
two protocols compared. (a) In MESI a remote read of a modified block forces a
writeback (red). (b) In MOESI the same event moves the holder to \stO\ without one
(green); the owner answers later reads and writes back only on eviction. All
other transitions are as in \cref{fig:mesi}.}
\label{fig:moesi-diff}
\end{figure}

\begin{table}[htbp]
\centering
\caption{Transition rules of the two protocols compared. Shaded rows are the only
ones that differ, so any measured difference between the runs comes from them.}
\label{tab:transition-diff}
\small
\begin{tabular}{@{}L{3.6cm} L{4.5cm} L{4.6cm}@{}}
\toprule
Event (state of other copies) & MESI & MOESI \\
\midrule
Read miss, no other copy & memory supplies; requester \stE & same \\
Read miss, copies in \stE/\stS & memory supplies; \stE$\to$\stS; requester \stS & same \\
\rowcolor{moesiOrange!12}
Read miss, a copy in \stM & holder supplies \textbf{and writes back}; holder \stS & holder supplies; holder \stO \\
\rowcolor{moesiOrange!12}
Read miss, a copy in \stO & (cannot occur) & owner supplies; owner stays \stO \\
Write hit in \stE / \stM & silent; requester \stM & same \\
Write hit in \stS (or \stO) & upgrade (8\,B) invalidates others; \stM & same \\
Write miss, a dirty copy exists & owner supplies; others invalidated; \stM & same \\
Write miss, no dirty copy & memory supplies; others invalidated; \stM & same \\
Evict \stS\ or \stE & silent & same \\
Evict \stM & write back (72\,B) & same \\
\rowcolor{moesiOrange!12}
Evict \stO & (cannot occur) & write back (72\,B) \\
\bottomrule
\end{tabular}
\end{table}

% ---------------------------------------------------------------
\section{What Is Held Fixed, and Paired Runs}
\label{sec:method-fixed}
\label{sec:method-paired}
\sectionmap{[Controlled variables [{Trace, caches, bus, latencies}] [Attribution test]
  [{One trace, two runs}] [Differences before averages]]}
\index{MESI}\index{MOESI}\index{Owned state}%

\Cref{tab:controlled} lists what is held constant and how the simulator of Part~V
enforces it. \Cref{fig:attribution-tree} turns the table into a test: a difference
between the two runs may be credited to \stO\ only if every question is answered
``yes''. The evaluations reviewed in Chapter~14 do not pass this test.

\begin{table}[htbp]
\centering
\caption{Controlled variables, identical in the MESI and MOESI runs of every
experiment.}
\label{tab:controlled}
\small
\begin{tabular}{@{}L{3.6cm} L{4.6cm} L{4.6cm}@{}}
\toprule
Variable & Value & Enforced by \\
\midrule
Memory-access trace & one list of (core, block, read/write) & generated once per seed, replayed twice \\
Order of accesses & trace order (traffic study); each core's own order (timing study) & one replay loop for both protocols \\
Cache per core & private, write-back, fully associative, LRU, 64-B blocks & one cache class; capacity a parameter \\
Interconnect & one snooping bus; broadcasts counted once & one bus model \\
Message sizes & 8\,B control, 72\,B data and writeback & constants of \cref{tab:counting-rules} \\
Latencies (timing study) & $T_{L1}$, $T_{\mathrm{mem}}$, $T_{\mathrm{coh}}$ of Chapters~2 and~8 & one parameter set per experiment \\
Accounting at the end & blocks still dirty are written back once & same drain step for both \\
\bottomrule
\end{tabular}
\end{table}

\begin{figure}[htbp]
\centering
\adjustbox{max width=0.95\linewidth}{%
\begin{forest}
  decision tree,
  [{Is the difference between\\the two runs caused by \stO?}
    [{Same trace and\\same interleaving?}
      [{Same caches, bus,\\message sizes, latencies?}, el=yes
        [{All rules not involving\\\stO\ identical?}, el=yes
          [{Both runs pass the\\coherence checks?}, el=yes
            [{Yes: attribute\\the difference to \stO}, el=yes]
            [{No: one run is wrong;\\fix before comparing}, el=no]]
          [{No: a second feature\\differs (not controlled)}, el=no]]
        [{No: configuration\\confounds the result}, el=no]]
      [{No: different sharing,\\not comparable}, el=no]]
  ]
\end{forest}}
\caption[Attribution test]{Decision tree for crediting a measured difference to
\stO. Each ``no'' leaf names the confounder that could explain the difference
instead.}
\label{fig:attribution-tree}
\end{figure}

Because one trace $W$ is replayed under both protocols, every configuration gives a
pair of results. For any quantity $X$ (bytes, writebacks, execution time) we
compute, per pair,
\begin{equation}
  \Delta X(W) = X(\mathrm{MESI},W) - X(\mathrm{MOESI},W),
  \qquad
  S_X(W) = \frac{\Delta X(W)}{X(\mathrm{MESI},W)},
  \label{eq:paired}
\end{equation}
and only then summarise over seeds by mean, minimum and maximum. For traffic,
$S_B$ is the fraction of MESI's bytes that MOESI saves; for time,
$T(\mathrm{MESI})/T(\mathrm{MOESI})$ is the speed-up. Traces differ from each other
far more than the two protocols differ on one trace, so differencing each pair
first keeps the trace-to-trace variation out of the comparison.

\label{sec:method-timing}
The traffic study replays one global access order, so both protocols see the same
interleaving. The timing study (Chapter~17) cannot do that and stay realistic. Each
core there has its own operation stream and clock, and the order in which accesses
reach the bus depends on how long earlier accesses took, as on real hardware.
MOESI makes some accesses faster, so its accesses can arrive at the bus in a
slightly different order. The per-core streams and think times are the same in both
runs; only the interleaving between cores may change. We therefore state and test
the exact results that need a common interleaving (\cref{prop:residency} and the
identities of Chapter~16) on the traffic study only, and report timing as paired
ratios over several seeds.

% ---------------------------------------------------------------
\section{The Two Runs Hold the Same Blocks}
\label{sec:method-residency}
\label{sec:method-example}
\sectionmap{[Same contents [Valid copies identical at every step]
  [Same misses and upgrades] [Only write backs differ] [Worked example]]}
\index{MESI}\index{MOESI}\index{Owned state}\index{Exclusive state}\index{forced writeback}\index{sharing patterns!migratory}\index{LRU replacement}\index{owner}\index{proof}%

The shaded rows of \cref{tab:transition-diff} might seem able to spread. If MOESI
kept a block that MESI dropped, the two runs would miss on different accesses and
the difference would no longer be due to \stO\ in any simple sense. This cannot
happen.

\begin{proposition}[Identical cache contents]
\label{prop:residency}
Replay the same trace under MESI and MOESI with the same cache model. After every
access, the set of (core, block) pairs held in a valid state (\stM, \stO, \stE\ or
\stS) is the same in both runs, and so is the LRU order of every cache. Both runs
therefore have exactly the same hits, read misses, write misses and upgrades.
\end{proposition}

\noindent The proof, by induction over the accesses, is in
\cref{app:proof-residency}. In short: the shaded read-miss rows install the block
in the requester and keep the supplier's copy valid in both protocols (as \stS\ in
MESI, \stO\ in MOESI), a write removes every other copy in both, and evictions are
chosen by the same LRU order. Under the counting rules of \cref{tab:counting-rules}
the two runs therefore send the same requests, upgrades and data messages and
differ only in writebacks. If $W_P$ is the number of writebacks of protocol $P$
(forced, eviction and end-of-run),
\begin{equation}
  \Delta B = B(\mathrm{MESI}) - B(\mathrm{MOESI}) = 72\,\bigl(W_{\mathrm{MESI}} - W_{\mathrm{MOESI}}\bigr)\ \text{bytes,}
  \label{eq:dB-writebacks}
\end{equation}
and the bytes written to memory differ by $64\,(W_{\mathrm{MESI}} -
W_{\mathrm{MOESI}})$. The proposition holds for the member of the class we chose. A
member that handed ownership to the reader or forwarded clean blocks would keep the
same contents, but one that dropped the owner's copy on a remote read would not.

\Cref{tab:method-example} applies the method by hand to a migratory block
(\cref{sec:dirty-patterns}): three cores read and then write it in turn, and P0
takes it a second time. The two runs make the same seven coherence transactions
(read misses at steps 1, 3, 5 and~7, upgrades at 4, 6 and~8). They differ only at
steps 3, 5 and~7, where MESI adds a 72-byte forced writeback, so the difference is
$632-416 = 216 = 3\times72$ bytes, as \cref{eq:dB-writebacks} predicts. Both pay
72 bytes at the end because P0, the last writer, still holds the block dirty and
nobody read it after step~8.

\begin{table}[htbp]
\centering
\caption{The method applied to a migratory block. The valid copies are the same in
both runs at every step (\cref{prop:residency}); only the labels \stS\ and \stO\
differ, and with them the bytes.}
\label{tab:method-example}
\small
\begin{tabular}{@{}c l c c c r r@{}}
\toprule
Step & Access & Valid copies & MESI states & MOESI states & MESI B & MOESI B \\
\midrule
1 & P0 reads  & P0     & E / I / I & E / I / I & 80  & 80 \\
2 & P0 writes & P0     & M / I / I & M / I / I & 0   & 0  \\
3 & P1 reads  & P0, P1 & S / S / I & \textbf{O} / S / I & 152 & 80 \\
4 & P1 writes & P1     & I / M / I & I / M / I & 8   & 8  \\
5 & P2 reads  & P1, P2 & I / S / S & I / \textbf{O} / S & 152 & 80 \\
6 & P2 writes & P2     & I / I / M & I / I / M & 8   & 8  \\
7 & P0 reads  & P0, P2 & S / I / S & S / I / \textbf{O} & 152 & 80 \\
8 & P0 writes & P0     & M / I / I & M / I / I & 8   & 8  \\
end & drain   & ---    & ---       & ---       & 72  & 72 \\
\midrule
 & \textbf{Total} & & & & \textbf{632} & \textbf{416} \\
\bottomrule
\end{tabular}
\end{table}

% ---------------------------------------------------------------
\section{Correctness Before Comparison}
\label{sec:method-correct}
\sectionmap{[Correctness [A wrong protocol can look fast]
  [Check the axioms on every access] [Version numbers] [Stop at the first violation]]}
\index{MESI}\index{MOESI}\index{MSI}\index{Owned state}\index{Exclusive state}\index{forced writeback}\index{Dirty Sharing Ratio}\index{coherence axioms}\index{owner}%

A protocol that skips a writeback it needs, or leaves a stale copy valid, looks
cheaper than a correct one, so a comparison means something only if both runs are
coherent. The checker tests this on every access instead of trusting the
transition table. After each access it looks at every copy of the block just
touched and confirms (1) at most one copy is in \stM, \stE\ or \stO; (2) a copy in
\stM\ or \stE\ is the only valid copy; and (3) MESI never holds \stO\ and MSI never
holds \stE\ or \stO. These are the single-writer, multiple-reader conditions that
the axioms of \S12.4.2 imply for an invalidation protocol~\cite{sarangi2025basic}.

State checks cannot catch a valid copy that holds old data. To check values
without simulating data, every block gets a version number. Each write increments
the block's latest version, each cache copy records the version it holds, and
memory records the version last written back. After every access the checker
confirms (4) every valid copy holds the latest version, so no read is stale, and
(5) if no cache holds the block in \stM\ or \stO, memory holds the latest version,
so no write is lost. Condition~5 is what the forced writeback protects in MESI and
what the owner protects in MOESI, and it catches the most likely MOESI bug, an
\stO\ copy evicted silently. \Cref{fig:versions} runs the checks on the six-step
trace of \cref{tab:trace}.

\begin{figure}[htbp]
\centering
\begin{bigfig}[0.95]
\begin{tikzpicture}[font=\sffamily\scriptsize,
    cell/.style={draw, minimum width=1.55cm, minimum height=0.5cm, align=center},
    hd/.style={font=\sffamily\bfseries\scriptsize}]
  \node[hd] at (-1.1,0.55) {after step};
  \foreach \i in {1,...,6} { \node[hd] at (\i*1.7,0.55) {\i}; }
  \node[hd, anchor=east] at (0.5,0) {latest};
  \node[hd, anchor=east, text=mesiBlue] at (0.5,-0.6) {MESI memory};
  \node[hd, anchor=east, text=moesiOrange] at (0.5,-1.2) {MOESI memory};
  \node[hd, anchor=east, text=moesiOrange] at (0.5,-1.8) {MOESI owner};
  \foreach \i/\l/\m/\o/\w in {1/v0/v0/v0/---, 2/v1/v0/v0/P0 (\stM), 3/v1/v1/v0/P0 (\stO), 4/v1/v1/v0/P0 (\stO), 5/v2/v1/v0/P1 (\stM), 6/v2/v2/v0/P1 (\stO)} {
    \node[cell, fill=coreGrey!8] at (\i*1.7,0) {\l};
    \node[cell, fill=mesiBlue!10] at (\i*1.7,-0.6) {\m};
    \node[cell, fill=moesiOrange!10] at (\i*1.7,-1.2) {\o};
    \node[cell, fill=moesiOrange!5] at (\i*1.7,-1.8) {\w};
  }
  \node[text=memGreen, align=left, anchor=west] at (0.1,-2.55)
    {MESI: memory is stale only while a copy is in \stM\ (steps 2 and 5), so condition 5 holds.\\
     MOESI: memory stays at v0, but an owner always holds the latest version, so condition 5 holds.};
\end{tikzpicture}
\end{bigfig}
\caption[Version-number check]{Versions on the six-step trace. Writes create new
versions at steps 2 and~5. In MESI the forced writebacks at steps 3 and~6 bring
memory up to date; in MOESI memory is never written during the trace, which is
correct only because an owner always has the latest version.}
\label{fig:versions}
\end{figure}

% ===============================================================
%  Chapter 16 -- The Dirty Sharing Ratio
% ===============================================================
\chapter{The Dirty Sharing Ratio: Predicting the Benefit of MOESI}
\label{ch:dsr}
\index{MOESI}\index{Dirty Sharing Ratio}%

\begin{figure}[htbp]
\centering
\adjustbox{max width=\linewidth}{%
\begin{forest}
  chapter map,
  [Chapter 16
    [Definition [One counter under MESI] [Worked example]]
    [Dirty epochs [Write to next write] [Where write backs fall]]
    [Saving identity [MOESI never sends more] [{$S=\sigma\cdot72\,\dsr/\bar b$}]]
    [Patterns {\&} mixtures [Survival fraction]]
    [From bytes to time [Two effects of \stO] [Bus load]]
    [Hypotheses [H1 to H4] [Our contributions]]
  ]
\end{forest}}
\caption[]{Structure of Chapter~16: the metric, its exact
meaning through dirty epochs, the identity linking it to MOESI's saving, the step
from bytes to time, and four testable hypotheses.}
\label{fig:ch16-map}
\end{figure}

% ---------------------------------------------------------------
\section{Definition}
\label{sec:dsr-def}
\sectionmap{[DSR [Coherence transactions under MESI] [Reads served by a modified copy]
  [Two counters] [Worked example]]}
\index{MESI}\index{MOESI}\index{forced writeback}\index{Dirty Sharing Ratio}\index{snooping protocol}%

RQ3 asks whether the benefit of removing forced writebacks can be predicted from a
workload's behaviour under MESI. A predictor that needs a MOESI run is no use, so
we want a quantity a MESI system can count with almost no extra hardware. Every bus
transaction under MESI is a read miss, a write miss or an upgrade, and we call
their total the number of coherence transactions $n_{\mathrm{tx}}$. Some read
misses find the block in another cache in \stM, and each of these causes a forced
writeback. We call their number $n_{\mathrm{dr}}$ (dirty reads).

\begin{keybox}[Definition: Dirty Sharing Ratio]
\[
  \dsr(W,N) \;=\; \frac{n_{\mathrm{dr}}}{n_{\mathrm{tx}}}
  \;=\; \frac{\text{read misses served by a cache holding the block in \stM}}
             {\text{read misses} + \text{write misses} + \text{upgrades}},
\]
measured while workload $W$ runs on $N$ cores under MESI.
\end{keybox}

On a snooping bus, a cache holding the requested block in \stM\ already announces
that it will supply and flush the data (\cref{fig:snoop-bus}). Counting the bus
cycles with that announcement gives $n_{\mathrm{dr}}$, and counting all
transactions gives $n_{\mathrm{tx}}$. Both events are visible to the bus arbiter,
so two counters are enough (\cref{fig:dsr-counter}). Our simulator keeps the two
counts directly. For the six-step trace priced in \cref{tab:mesi-bill}, MESI makes
five coherence transactions (read misses at steps 1, 3, 4 and~6 and the upgrade at
step~5), and two of the read misses (steps 3 and~6) find the block in \stM, so
\begin{equation}
  \dsr = \tfrac{2}{5} = 0.40 .
  \label{eq:dsr-example}
\end{equation}

\begin{figure}[htbp]
\centering
\begin{bigfig}[0.92]
\begin{tikzpicture}[font=\sffamily\scriptsize,
    cc/.style={draw, thick, rounded corners=2pt, minimum width=1.6cm, minimum height=0.8cm, align=center},
    ar/.style={-{Stealth}, thick}]
  \draw[line width=3pt, coreGrey!50] (-0.6,0) -- (11.8,0);
  \node[font=\sffamily\bfseries\scriptsize, text=coreGrey] at (12.4,0) {bus};
  \node[cc, fill=tealC!12, draw=tealC] (r) at (0.8,1.5) {P1\\requester};
  \node[cc, fill=mesiBlue!12, draw=mesiBlue] (h) at (3.6,1.5) {P0\\holds \stM};
  \node[cc, fill=coreGrey!8, draw=coreGrey] (p2) at (6.4,1.5) {P2\\---};
  \node[cc, fill=amberC!14, draw=amberC] (mem) at (9.4,1.5) {Memory};
  \foreach \n in {r,h,p2,mem} { \draw[thick, coreGrey!70] (\n.south) -- (\n.south|-0,0); }
  \draw[ar, tealC] (r.south) ++(0.3,0) -- ++(0,-0.95) node[below, font=\sffamily\tiny]{readmiss};
  \draw[ar, alertRed] (h.south) ++(0.3,0) -- ++(0,-0.95) node[below, font=\sffamily\tiny, align=center]{``dirty''\\asserted};
  \node[cc, fill=memGreen!12, draw=memGreen, minimum width=2.6cm] (cnt) at (4.8,-2.3)
     {counters in the arbiter\\$n_{\mathrm{tx}}$ \texttt{+= 1} on every transaction\\$n_{\mathrm{dr}}$ \texttt{+= 1} when ``dirty'' asserted};
  \draw[ar, memGreen] (3.9,-1.25) -- (cnt.north west);
  \draw[ar, memGreen] (1.1,-1.0) to[out=-90, in=180] (cnt.west);
  \node[text=memGreen, font=\sffamily\bfseries\small] at (9.6,-2.3) {$\dsr = n_{\mathrm{dr}} / n_{\mathrm{tx}}$};
\end{tikzpicture}
\end{bigfig}
\caption[Measuring the Dirty Sharing Ratio]{Measuring DSR on a MESI snooping bus.
The holder of a modified block already signals that it will supply the data, so
two counters give the metric.}
\label{fig:dsr-counter}
\end{figure}

\Cref{tab:dsr-example} carries the example through to MOESI. MOESI saves 72 of
MESI's 472 bytes, or 15.3\,\%, which is half the 30.5\,\% of forced writebacks
counted in Chapter~7. The forced writeback at step~3 disappears, because P0's owned
copy is made obsolete by P1's write at step~5 and never needs to reach memory. The
one at step~6 is only deferred: P1 still owns the block when the run ends and pays
the writeback then.

\begin{table}[htbp]
\centering
\caption{The six-step trace of \cref{tab:trace} under MESI and MOESI, with bytes
from \cref{tab:counting-rules}. ``Drain'' charges one writeback at the end for
each block still dirty.}
\label{tab:dsr-example}
\small
\begin{tabular}{@{}c l c c c c l@{}}
\toprule
Step & Access & Transaction & Dirty read? & MESI B & MOESI B & Epoch \\
\midrule
1 & P0 reads  & read miss & no  & 80  & 80 & --- \\
2 & P0 writes & (silent)  & --- & 0   & 0  & 1 opens (P0) \\
3 & P1 reads  & read miss & \textbf{yes} & 152 & 80 & 1 \\
4 & P2 reads  & read miss & no  & 80  & 80 & 1 \\
5 & P1 writes & upgrade   & --- & 8   & 8  & 1 ends; 2 opens (P1) \\
6 & P0 reads  & read miss & \textbf{yes} & 152 & 80 & 2 \\
end & drain   & ---       & --- & 0   & 72 & 2 ends \\
\midrule
 & \textbf{Total} & 5 & 2 & \textbf{472} & \textbf{400} & \\
\bottomrule
\end{tabular}
\end{table}

% ---------------------------------------------------------------
\section{Dirty Epochs}
\label{sec:dsr-epochs}
\sectionmap{[Epochs [From one write to the next] [At most one write back each]
  [MESI pays if anyone reads] [MOESI pays if the owner leaves]]}
\index{MESI}\index{MOESI}\index{Owned state}\index{forced writeback}\index{dirty epoch}\index{owner}%

Fix one block. A \gls{epoch}\index{dirty epoch} starts when some core writes the
block and ends at the next write to it by any core, or at the end of the run.
During an epoch the value does not change, and one core, the last writer, holds the
newest data when the epoch starts. Epochs divide the whole history of a block, so
each writeback of the block falls in some epoch (\cref{fig:epoch}). Both protocols
write a block back only when a dirty copy stops being dirty, and after the first
writeback in an epoch no dirty copy is left until the next write, so each protocol
makes at most one writeback per epoch. MESI writes back if the epoch has at least
one remote read (the forced writeback, at the first one), or otherwise if the
writer evicts its \stM\ copy before the epoch ends. MOESI writes back if and only if
the owner's copy (\stM\ or \stO) is evicted, or the run ends, before the epoch ends.
The two differ in one case only: an epoch with at least one remote read in which
the owner keeps the block until the next write. MESI pays a writeback there and
MOESI does not, and we call such an epoch \emph{saved}.

\begin{figure}[htbp]
\centering
\begin{bigfig}[0.92]
\begin{tikzpicture}[font=\sffamily\scriptsize]
  \draw[-{Stealth}, thick, coreGrey] (0,0) -- (12.6,0) node[right]{time};
  \fill[memGreen!10] (0.6,-0.35) rectangle (11.4,0.35);
  \node[font=\sffamily\bfseries\scriptsize, text=memGreen] at (6.0,0.62) {one dirty epoch (value unchanged)};
  \foreach \x/\t in {0.6/{P0 writes}, 11.4/{P2 writes}} {
     \draw[very thick, alertRed] (\x,-0.4) -- (\x,0.4); \node[text=alertRed, below] at (\x,-0.4) {\t}; }
  \foreach \x/\t in {3.0/{P1 reads}, 5.6/{P2 reads}, 8.2/{P3 reads}} {
     \draw[thick, tealC] (\x,-0.3) -- (\x,0.3); \node[text=tealC, below] at (\x,-0.35) {\t}; }
  \node[anchor=east, text=mesiBlue, font=\sffamily\bfseries\scriptsize] at (0.3,-1.5) {MESI};
  \draw[mesiBlue, thick] (0.6,-1.5) -- (11.4,-1.5);
  \node[draw=mesiBlue, fill=mesiBlue!15, rounded corners=1pt, inner sep=2pt] at (3.0,-1.5) {forced WB};
  \node[text=mesiBlue, align=center] at (5.6,-2.05) {memory\\supplies};
  \node[text=mesiBlue, align=center] at (8.2,-2.05) {memory\\supplies};
  \node[text=mesiBlue, anchor=west] at (0.6,-2.65) {1 writeback, at the first remote read, whatever happens later};
  \node[anchor=east, text=moesiOrange, font=\sffamily\bfseries\scriptsize] at (0.3,-3.5) {MOESI};
  \draw[moesiOrange, thick] (0.6,-3.5) -- (11.4,-3.5);
  \foreach \x in {3.0,5.6,8.2} { \node[text=moesiOrange, align=center] at (\x,-4.05) {owner\\supplies}; }
  \node[draw=moesiOrange, fill=moesiOrange!15, rounded corners=1pt, inner sep=2pt] at (11.4,-3.5) {invalidated};
  \node[text=moesiOrange, anchor=west] at (0.6,-4.65) {0 writebacks if the owner keeps the block until the next write; 1 if it is evicted first};
\end{tikzpicture}
\end{bigfig}
\caption[A dirty epoch]{One dirty epoch with three remote reads. MESI pays one
writeback at the first remote read and memory serves later readers. MOESI pays
nothing while the owner keeps the block, and the value becomes obsolete at the next
write.}
\label{fig:epoch}
\end{figure}

% ---------------------------------------------------------------
\section{The Saving Identity}
\label{sec:dsr-identity}
\sectionmap{[Identity [MOESI never sends more bytes] [Saving = 72 per saved epoch]
  [Survival fraction sigma] [{$S=\sigma\cdot72\,\dsr/\bar b$}]]}
\index{MESI}\index{MOESI}\index{Dirty Sharing Ratio}\index{survival fraction}\index{dirty epoch}\index{sharing patterns!producer--consumer}\index{sharing patterns!migratory}\index{owner}\index{proof}%

\begin{proposition}[MOESI never sends more bytes]
\label{prop:never-worse}
Under the method of Chapter~15 and the counting rules of
\cref{tab:counting-rules}, for every trace $W$,
\[
  0 \;\le\; \Delta B(W) \;=\; 72\, n_{\mathrm{saved}}(W) \;\le\; 72\, n_{\mathrm{dr}}(W),
\]
where $n_{\mathrm{saved}}$ is the number of saved epochs and $n_{\mathrm{dr}}$ the
number of dirty reads under MESI.
\end{proposition}

\noindent The proof is in \cref{app:proof-never-worse}. It compares the two
protocols epoch by epoch using \cref{prop:residency}, and notes that each saved
epoch contains exactly one dirty read. Two consequences follow. In the traffic
model MOESI cannot lose, so any loss has to come from timing (\cref{sec:dsr-time}).
And the dirty reads that DSR counts are exactly the epochs MOESI could save. The
fraction it does save is the \gls{sigma}
\begin{equation}
  \sigma = \frac{n_{\mathrm{saved}}}{n_{\mathrm{dr}}} \in [0,1],
  \label{eq:sigma}
\end{equation}
the share of dirty-read epochs whose owner keeps the block until the next write.
Large caches and frequent rewriting push $\sigma$ towards~1, while small caches,
which evict owners early, and runs that end with many blocks still owned push it
down. Dividing $\Delta B$ by MESI's total bytes, writing $\bar b =
B_{\mathrm{MESI}}/n_{\mathrm{tx}}$ for the average bytes MESI moves per coherence
transaction (writebacks included), and using $n_{\mathrm{dr}} = \dsr\cdot
n_{\mathrm{tx}}$ gives the central result of this chapter.

\begin{keybox}[The saving identity]
\begin{equation}
  S_B \;=\; \frac{\Delta B}{B_{\mathrm{MESI}}}
      \;=\; \sigma\,\frac{72\; n_{\mathrm{dr}}}{B_{\mathrm{MESI}}}
      \;=\; \sigma \cdot F \;=\; \sigma\cdot\frac{72\,\dsr}{\bar b}.
  \label{eq:identity}
\end{equation}
The share of traffic MOESI saves equals the survival fraction times the
forced-writeback fraction of \cref{eq:fwb-fraction}. Equivalently, it is
proportional to DSR with constant $72\sigma/\bar b$.
\end{keybox}

The identity agrees with the earlier examples. For the six-step trace, $\dsr=0.4$,
$\bar b = 472/5 = 94.4$\,B and $\sigma = 1/2$ (one of the two dirty-read epochs is
saved), so $0.5\times72\times0.4/94.4 = 15.3\,\%$, as in \cref{tab:dsr-example}.
For the migratory trace of \cref{tab:method-example} the block has four epochs,
opened at steps 2, 4, 6 and~8. The first three each have one remote read and end at
the next write with the owner still holding the block, so all three are saved; the
fourth has no remote read. Then $n_{\mathrm{dr}}=3$, $n_{\mathrm{tx}}=7$, $\sigma=1$
and $\bar b = 632/7 \approx 90.3$\,B, giving $72\times(3/7)/90.3 = 216/632 =
34.2\,\%$. For the producer--consumer model of \cref{eq:pc-model} with $C$
consumers, each round has $C+1$ transactions of which one is a dirty read, so
$\dsr = 1/(C+1)$ and $\bar b = (160+80(C-1))/(C+1)$. In steady state the producer
always keeps the block, so $\sigma=1$ and \cref{eq:identity} becomes
$72/(160+80(C-1))$, which is 45\,\% for $C=1$, as in \cref{tab:traffic-patterns}.

\Cref{eq:identity} is exact, but besides DSR it contains $\bar b$, which is easy to
measure under MESI, and $\sigma$, which depends on what happens after each dirty
read. The predictor we propose is therefore
\begin{equation}
  \hat S_B = \frac{72\,\dsr}{\bar b}\qquad(\sigma\text{ taken as }1),
  \label{eq:predictor}
\end{equation}
which overestimates the saving by the factor $1/\sigma$. When $\bar b$ varies little
between workloads, $\hat S_B$ is close to a straight line in DSR with slope
$72/\bar b$ (\cref{fig:identity-lines}). How close $\sigma$ stays to~1 and how much
$\bar b$ varies cannot be derived in general; Part~VI measures both.

\begin{figure}[htbp]
\centering
\adjustbox{max width=\linewidth}{%
\begin{tikzpicture}
  \begin{axis}[width=10.5cm, height=5.4cm,
      xlabel={Dirty Sharing Ratio under MESI}, ylabel={Predicted saving $\hat S_B$ (\%)},
      xmin=0, xmax=0.5, ymin=0, ymax=46, xtick={0,0.1,0.2,0.3,0.4,0.5},
      xticklabel style={/pgf/number format/fixed, /pgf/number format/precision=1},
      legend style={at={(1.02,0.5)}, anchor=west, font=\scriptsize, draw=none},
      grid=major, grid style={coreGrey!15},
      label style={font=\small}, tick label style={font=\scriptsize}]
    \addplot[very thick, tealC, domain=0:0.5] {7200*x/80};   \addlegendentry{$\bar b = 80$\,B}
    \addplot[very thick, plumC, domain=0:0.5] {7200*x/100};  \addlegendentry{$\bar b = 100$\,B}
    \addplot[very thick, amberC, domain=0:0.5] {7200*x/130}; \addlegendentry{$\bar b = 130$\,B}
    \addplot[only marks, mark=*, mark size=2.4pt, moesiOrange] coordinates
      {(0.5,45) (0.3333,30) (0.2,18) (0.1111,10) (0.0625,5.625)};
    \addlegendentry{prod.--cons., $C=1,2,4,8,15$}
    \addplot[only marks, mark=square*, mark size=2.6pt, mesiBlue] coordinates {(0.4,30.5)};
    \addlegendentry{six-step trace, $\sigma=1$}
  \end{axis}
\end{tikzpicture}}
\caption[The predictor as a family of lines]{\Cref{eq:predictor} for three values of
$\bar b$. The producer--consumer model lies on the 80\,B line because each round
moves $80(C+1)$ bytes in $C+1$ transactions. The blue square is the six-step trace
($\bar b=94.4$\,B) if its last epoch were also saved. Analytical.}
\label{fig:identity-lines}
\end{figure}

% ---------------------------------------------------------------
\section{The Sharing Patterns and Their Mixtures}
\label{sec:dsr-patterns}
\sectionmap{[Patterns [DSR of each pattern] [Mixtures weighted by transactions]
  [A little dirty sharing goes a long way] [Survival fraction]]}
\index{MESI}\index{MOESI}\index{dirty sharing}\index{Dirty Sharing Ratio}\index{survival fraction}\index{dirty epoch}\index{sharing patterns!producer--consumer}\index{sharing patterns!migratory}\index{LRU replacement}\index{owner}%

\Cref{tab:dsr-patterns} applies the definition to one unit of each sharing pattern
of \cref{tab:patterns}, in steady state and with caches large enough that owners
are not evicted; the entries follow from the transaction counts of
\cref{tab:traffic-patterns}. Private and read-only data never involve dirty data
read by another core, so their DSR is zero and MOESI saves nothing on them. The two
dirty-sharing patterns both have $\bar b = 80$\,B, which is why their savings lie on
one line of \cref{fig:identity-lines}. Migratory sharing has the largest DSR in the
table, one half, since every hand-off is one dirty read and one upgrade.

\begin{table}[htbp]
\centering
\caption{The Dirty Sharing Ratio and the saving identity for one unit of each
sharing pattern ($C$ consumers per producer write), in steady state with no
capacity evictions. $S_B$ from \cref{eq:identity}.}
\label{tab:dsr-patterns}
\small
\begin{tabular}{@{}L{3.1cm} c c c c c c@{}}
\toprule
Pattern (unit) & $n_{\mathrm{tx}}$ & $n_{\mathrm{dr}}$ & $\dsr$ & $\bar b$ (B) &
  $\sigma$ & $S_B$ \\
\midrule
Private (read, write, evict) & 1 & 0 & 0 & 152 & --- & 0 \\
Read-only (per reader) & 1 & 0 & 0 & 80 & --- & 0 \\
Producer--consumer (per round) & $C+1$ & 1 & $\frac{1}{C+1}$ & 80 & 1 &
  $\frac{0.9}{C+1}$ \\
Migratory (per hand-off) & 2 & 1 & $\frac{1}{2}$ & 80 & 1 & 45\,\% \\
\bottomrule
\end{tabular}
\end{table}

A real program mixes the patterns. DSR is a ratio of transaction counts, so the DSR
of a mixture is the average of the patterns' DSRs weighted by the transactions each
produces, not by its accesses:
\begin{equation}
  \dsr_{\mathrm{mix}} = \frac{\sum_i n_{\mathrm{dr},i}}{\sum_i n_{\mathrm{tx},i}}
  = \frac{\sum_i n_{\mathrm{tx},i}\,\dsr_i}{\sum_i n_{\mathrm{tx},i}} .
  \label{eq:dsr-mix}
\end{equation}
Private data that mostly hits produces few transactions, so a small amount of
migratory or producer--consumer sharing can dominate the DSR. If a fraction $m$ of
all accesses go to migratory blocks (one transaction each, $\dsr = 1/2$) and the
rest are private accesses with hit rate $h$ ($1-h$ transactions each, $\dsr=0$),
\begin{equation}
  \dsr_{\mathrm{mix}}(m,h) = \frac{m/2}{m + (1-m)(1-h)} .
  \label{eq:dsr-mix-mig}
\end{equation}
With a 99\,\% private hit rate, 1\,\% of migratory accesses already gives a DSR of
0.25, half the maximum (\cref{fig:dsr-mix}). A workload's DSR can therefore be high
even when sharing is a small part of its accesses, and the identity ties MOESI's
saving to the DSR, not to the share of accesses.

\begin{figure}[htbp]
\centering
\begin{minipage}[c]{0.54\linewidth}\centering
\begin{tikzpicture}
  \begin{axis}[width=\linewidth, height=5.0cm,
      xlabel={Accesses to migratory blocks, $m$ (\%)}, ylabel={$\dsr_{\mathrm{mix}}$},
      xmin=0, xmax=20, ymin=0, ymax=0.6,
      legend style={at={(0.98,0.04)}, anchor=south east, font=\tiny, draw=none, fill=white},
      grid=major, grid style={coreGrey!15},
      label style={font=\footnotesize}, tick label style={font=\scriptsize}]
    \addplot[very thick, tealC, domain=0:20, samples=120] {(x/100/2)/((x/100) + (1-x/100)*0.10)};
      \addlegendentry{private hit rate $h = 90\,\%$}
    \addplot[very thick, plumC, domain=0:20, samples=120] {(x/100/2)/((x/100) + (1-x/100)*0.03)};
      \addlegendentry{$h = 97\,\%$}
    \addplot[very thick, brickC, domain=0:20, samples=160] {(x/100/2)/((x/100) + (1-x/100)*0.01)};
      \addlegendentry{$h = 99\,\%$}
    \addplot[thick, dashed, coreGrey, forget plot] coordinates {(0,0.5) (20,0.5)};
  \end{axis}
\end{tikzpicture}
\caption[DSR of a mixture]{\Cref{eq:dsr-mix-mig}: a few per cent of migratory
accesses dominate the metric (dashed: the migratory maximum, 0.5). Analytical.}
\label{fig:dsr-mix}
\end{minipage}\hfill
\begin{minipage}[c]{0.43\linewidth}\centering
\adjustbox{max width=\linewidth}{%
\begin{forest}
  decision tree,
  [{One dirty epoch}
    [{Any remote read\\before the next write?}
      [{Owner evicted\\before next write?}, el=yes
        [{Deferred\\(lowers $\sigma$)}, el=yes]
        [{Run ends\\first?}, el=no
          [{Deferred:\\paid at end}, el=yes]
          [{Saved}, el=no]]]
      [{Nothing\\to save}, el=no]]
  ]
\end{forest}}
\caption[Is an epoch saved?]{Classifying a dirty epoch. The fraction of dirty-read
epochs that are saved is $\sigma$.}
\label{fig:sigma-tree}
\end{minipage}
\end{figure}

The remaining unknown is $\sigma$. \Cref{fig:sigma-tree} gives the test for one
epoch. Under the method of Chapter~15 the owner's copy is evicted at the same moment
in both protocols (\cref{prop:residency}), so $\sigma$ is a property of the workload
and the cache, not of the protocol. Two things lower it. The first is capacity: if
the owner touches more distinct blocks during the epoch than its cache can hold, LRU
evicts the owned block before the next write, and the writeback MESI paid early is
paid late instead. Small caches, large working sets and long epochs all make this
more likely. The second is the end of the run: epochs still open when measurement
stops are charged a writeback, as at step~6 of \cref{tab:dsr-example}. Over a long
run this is a small fixed number of epochs, so its effect shrinks, but in a short
trace such as the six-step one it can halve the saving.

% ---------------------------------------------------------------
\section{From Bytes to Time}
\label{sec:dsr-time}
\sectionmap{[Time [O changes two things] [Write back removed]
  [Owner answers later reads] [Latency per epoch] [Load decides]]}
\index{MESI}\index{MOESI}\index{Owned state}\index{forced writeback}\index{dirty epoch}\index{AMAT}\index{bus utilisation}\index{cache-to-cache transfer}\index{Effect A and Effect B}\index{coherence miss}%

RQ2 asks whether saved bytes become saved time, and RQ4 whether deferring
writebacks can cost more than it saves. With the controlled method we can see that
\stO\ changes two things, with different signs (\cref{fig:two-effects}).

\begin{figure}[htbp]
\centering
\adjustbox{max width=0.95\linewidth}{%
\begin{forest}
  chapter map,
  [What \stO\ changes
    [{Effect A: write back removed}
      [{at the first remote read of an epoch}]
      [{saves bus and memory occupancy}]
      [{and the reader's wait for the flush}]
      [{sign: always helps}]]
    [{Effect B: owner answers later reads}
      [{readers 2, 3, ... of the epoch}]
      [{cache-to-cache instead of memory}]
      [{faster if $T_{\mathrm{c2c}} < T_{\mathrm{mem}}$}]
      [{sign: helps or hurts}]]
  ]
\end{forest}}
\caption[Two effects of the Owned state]{The \stO\ state bundles two effects.
Effect~A, removing the forced writeback, is the subject of this report; Effect~B
changes the supplier for later readers. Only Effect~A always helps.}
\label{fig:two-effects}
\end{figure}

Take a saved epoch with $r\ge1$ remote reads on an otherwise idle bus. Let
$T_{\mathrm{coh}}$ be the time for a read served by a cache holding the block in
\stM\ under MESI, including the writeback (as in \cref{eq:amat-coh}); let
$t_{\mathrm{wb}}$ be the part spent transferring the writeback; let
$T_{\mathrm{c2c}} = T_{\mathrm{coh}} - t_{\mathrm{wb}}$ be a cache-to-cache read
without a writeback; and let $T_{\mathrm{mem}}$ be a read served by memory. Under
MESI the first reader waits $T_{\mathrm{coh}}$ and the other $r-1$ wait
$T_{\mathrm{mem}}$, while under MOESI all $r$ wait $T_{\mathrm{c2c}}$. The reader
latency saved in the epoch is
\begin{equation}
  \Delta L(r) \;=\; \underbrace{\bigl(T_{\mathrm{coh}} - T_{\mathrm{c2c}}\bigr)}_{\text{Effect A}\;=\;t_{\mathrm{wb}}}
  \;+\; \underbrace{(r-1)\bigl(T_{\mathrm{mem}} - T_{\mathrm{c2c}}\bigr)}_{\text{Effect B}} .
  \label{eq:dL}
\end{equation}
With the assumed $T_{\mathrm{coh}}=60$ and $T_{\mathrm{mem}}=100$ cycles (Chapters~2
and~8) and the 9-cycle transfer of a 72-byte writeback on the 8-byte-per-cycle bus
of Chapter~17, Effect~A is 9 cycles per saved epoch and Effect~B is $100-51 = 49$
cycles per extra reader. Effect~B changes sign once a cache-to-cache read is slower
than a memory read (\cref{fig:dL}). With $T_{\mathrm{c2c}}=141$ cycles, for
example, an epoch with three readers loses $9-2\times41 = -73$ cycles under MOESI.

\begin{figure}[htbp]
\centering
\adjustbox{max width=\linewidth}{%
\begin{tikzpicture}
  \begin{axis}[width=10cm, height=5.2cm,
      xlabel={Remote reads per saved epoch, $r$}, ylabel={$\Delta L(r)$ (cycles)},
      xmin=1, xmax=8, ymin=-320, ymax=360, xtick={1,...,8},
      legend style={at={(1.02,0.5)}, anchor=west, font=\scriptsize, draw=none},
      grid=major, grid style={coreGrey!15},
      label style={font=\small}, tick label style={font=\scriptsize}]
    \addplot[fill=alertRed!8, draw=none, forget plot] coordinates {(1,0) (8,0) (8,-320) (1,-320)} \closedcycle;
    \addplot[very thick, memGreen, mark=*, domain=1:8, samples=8] {9 + (x-1)*(100-51)};
      \addlegendentry{$T_{\mathrm{c2c}} = 51$}
    \addplot[very thick, amberC, mark=square*, domain=1:8, samples=8] {9 + (x-1)*(100-91)};
      \addlegendentry{$T_{\mathrm{c2c}} = 91$}
    \addplot[very thick, alertRed, mark=triangle*, domain=1:8, samples=8] {9 + (x-1)*(100-141)};
      \addlegendentry{$T_{\mathrm{c2c}} = 141$}
    \addplot[thick, dashed, coreGrey, forget plot] coordinates {(1,0) (8,0)};
  \end{axis}
\end{tikzpicture}}
\caption[Latency saved per epoch]{\Cref{eq:dL} with $T_{\mathrm{mem}}=100$ and
$t_{\mathrm{wb}}=9$ cycles. Effect~A gives each saved epoch 9 cycles; Effect~B adds
or subtracts $T_{\mathrm{mem}}-T_{\mathrm{c2c}}$ per further reader (shaded: MOESI
slower). Analytical, idle bus.}
\label{fig:dL}
\end{figure}

\Cref{eq:dL} ignores waiting. On a bus that carries one message at a time
(\cref{eq:util}) every transaction also occupies the bus, and under MESI a dirty
read occupies it for one more 72-byte message. On a lightly used bus requests rarely
wait, so the occupancy MOESI saves shortens almost nothing, and execution time
changes only by the latency of \cref{eq:dL}, spread over a run dominated by
computation and hits. As utilisation nears 100\,\% waiting grows sharply
(Chapter~8), execution time is set by the total time the bus is busy, and the
speed-up of MOESI tends to the ratio of bus-busy times,
\begin{equation}
  \frac{T(\mathrm{MESI})}{T(\mathrm{MOESI})} \;\longrightarrow\;
  \frac{C_{\mathrm{bus}}(\mathrm{MESI})}{C_{\mathrm{bus}}(\mathrm{MOESI})}
  \qquad\text{as } U \to 100\,\%,
  \label{eq:sat-limit}
\end{equation}
where $C_{\mathrm{bus}}$ is the total number of cycles the bus is occupied. By
\cref{prop:residency} the two runs differ in $C_{\mathrm{bus}}$ only by the
occupancy of the writebacks MOESI avoids. The same holds for the memory port, since
every forced writeback is also a 64-byte memory write.

\Cref{eq:amat-coh} shows the low-load case in numbers. Chapter~8 used the 3-cycle
single-core AMAT, $T_{\mathrm{coh}}=60$ cycles and a coherence-miss rate of 1, 2, 4
and 8\,\% at 2, 4, 8 and 16 cores (\cref{fig:amat-bars}). In the case most
favourable to Effect~A, every coherence miss is a dirty read with one reader
($r=1$), so under MOESI each costs $T_{\mathrm{c2c}}=51$ cycles instead of 60
(\cref{tab:amat-compare}). Even then, at 16 cores Effect~A saves 0.72 of 7.8 cycles,
under 10\,\%. A larger time gain has to come from Effect~B or from bus load
(\cref{eq:sat-limit}), which is why H2 below ties the gain to utilisation.

\begin{table}[htbp]
\centering
\caption{AMAT from \cref{eq:amat-coh} with the illustrative values of Chapter~8,
when every coherence miss is a dirty read with one reader. Idle interconnect;
cycles.}
\label{tab:amat-compare}
\small
\begin{tabular}{@{}c c c c c c@{}}
\toprule
Cores & $m_{\mathrm{coh}}$ & AMAT, MESI & AMAT, MOESI & Saved (cycles) & Saved (\%) \\
\midrule
2  & 1\,\% & 3.60 & 3.51 & 0.09 & 2.5 \\
4  & 2\,\% & 4.20 & 4.02 & 0.18 & 4.3 \\
8  & 4\,\% & 5.40 & 5.04 & 0.36 & 6.7 \\
16 & 8\,\% & 7.80 & 7.08 & 0.72 & 9.2 \\
\bottomrule
\end{tabular}
\end{table}

% ---------------------------------------------------------------
\section{Hypotheses and Contributions}
\label{sec:dsr-hypotheses}
\label{sec:contrib-summary}
\sectionmap{[Hypotheses [H1: F falls with cores] [H2: load gates the time gain]
  [H3: DSR predicts the saving] [H4: MOESI loses only through timing] [Contributions C1--C16]]}
\index{MOESI}\index{Dirty Sharing Ratio}\index{survival fraction}\index{directory protocol}\index{sharing patterns!migratory}\index{cache-to-cache transfer}%

The analysis turns each research question of \cref{tab:rq} into a hypothesis that
Part~VI tests (\cref{tab:hypotheses}). We state them here, before any measurement,
so that Part~VI can report where they fail.

\begin{table}[htbp]
\centering
\caption{Hypotheses derived in Part~IV, one per research question, and how each is
tested in Part~VI.}
\label{tab:hypotheses}
\small
\begin{tabular}{@{}l L{5.6cm} L{2.6cm} L{3.4cm}@{}}
\toprule
 & Hypothesis & Basis & Test in Part VI \\
\midrule
H1 (RQ1) & For a fixed sharing structure, $F(W,N)$ falls as $N$ grows,
  because more readers share each epoch: one forced writeback, more 80-byte
  read misses. & \cref{eq:pc-model}, \cref{sec:dsr-epochs} & $F$ measured for
  $N = 2$ to 16 \\
H2 (RQ2) & The execution-time gain of MOESI is near zero at low bus
  utilisation and rises towards the bus-occupancy ratio as utilisation
  approaches 100\,\%. & \cref{eq:dL,eq:sat-limit} & load sweep: gain against
  utilisation \\
H3 (RQ3) & The traffic saving obeys \cref{eq:identity} exactly; with
  $\sigma$ near 1 and $\bar b$ nearly constant, DSR predicts it through
  \cref{eq:predictor}. & \cref{prop:never-worse}, \cref{eq:identity} & DSR
  against measured $S_B$; error of $\hat S_B$ \\
H4 (RQ4) & MOESI never moves more bytes than MESI; it can take more time
  only through Effect~B (when $T_{\mathrm{c2c}} > T_{\mathrm{mem}}$) or owner
  contention. Small caches reduce the gain by lowering $\sigma$ but cannot
  reverse it in bytes. & \cref{prop:never-worse}, \cref{eq:dL} & cache-size
  sweep; latency sensitivity sweep \\
\bottomrule
\end{tabular}
\end{table}

DSR is not the first measure of sharing. Archibald and Baer parameterised their
simulations by the fraction of references to shared
blocks~\cite[p.~286]{archibald1986evaluation}, and adaptive protocols detect
migratory blocks at run time, for example by noticing a write by a processor other
than the last writer~\cite{stenstrom1993adaptive,nilsson1994adaptive}.
\Cref{tab:dsr-vs} shows how DSR differs from these and from $F$. The identity
\cref{eq:identity} is exact only under the method of Chapter~15: one MOESI member,
the same cache contents in both runs, and the counting rules of
\cref{tab:counting-rules}. A directory protocol, another MOESI member or a
non-inclusive hierarchy would change $\bar b$ and could change $\sigma$. DSR also
predicts traffic only; for time, bus load and the relative speed of cache-to-cache
and memory transfers matter as well (\cref{sec:dsr-time}). Part~VII returns to
these limits.

\begin{table}[htbp]
\centering
\caption{The Dirty Sharing Ratio compared with related measures.}
\label{tab:dsr-vs}
\small
\begin{tabular}{@{}L{3.1cm} L{3.0cm} L{3.4cm} L{3.4cm}@{}}
\toprule
Measure & What it counts & Used for & Limitation for RQ3 \\
\midrule
Shared-reference fraction~\cite{archibald1986evaluation} & references to
  shared blocks & input parameter of a model & ignores whether shared data is
  dirty when read \\
Migratory detection~\cite{stenstrom1993adaptive,nilsson1994adaptive} & blocks
  written by successive processors & run-time protocol adaptation & per block,
  inside one protocol; not a workload score \\
Forced-writeback fraction $F$ (\cref{eq:fwb-fraction}) & bytes of forced
  writebacks & definition of the problem & needs byte accounting of all
  messages \\
Dirty Sharing Ratio & dirty reads per transaction & choosing between
  protocols before building one & needs $\bar b$; assumes $\sigma \approx 1$ \\
\bottomrule
\end{tabular}
\end{table}

\label{sec:contrib-parts123}
\label{sec:contrib-parts45}
\Cref{tab:contrib-list} collects what this report adds to the sources it reviews.
C1 to C10 are analytical results reached while studying existing work in Parts~I
to~III; the material they are built from belongs to the cited sources, and the
results are ours. C11 to C16 are the method, metric and tool of Parts~IV and~V. All
analytical results rest on the message sizes of \cref{tab:msg-classes} and the
counting rules of \cref{tab:counting-rules}, including our assumed 16-byte word
update, and C3 to C10 ignore contention, latency and capacity, so they say nothing
about execution time on their own.

\begin{table}[htbp]
\centering
\caption{Contributions of this report. ``Basis'' names the published material
each result is built from.}
\label{tab:contrib-list}
\label{tab:contrib-45}
\footnotesize
\renewcommand{\arraystretch}{1.12}
\begin{tabular}{@{}l L{4.1cm} L{4.6cm} L{3.4cm}@{}}
\toprule
 & Result & What it shows & Where; basis \\
\midrule
C1 & Full MESI state diagrams & completes the textbook's ``exercise for the
  reader'' & \cref{fig:mesi}; \cite[p.~603]{sarangi2025basic} \\
C2 & One trace under MSI, MESI and MOESI & the protocols differ at three points:
  \stE, silent upgrade, \stO & \cref{tab:trace}; \cite{papamarcos1984illinois,sweazey1986moesi} \\
C3 & The trace priced under MESI & writebacks are 30.5\,\% of 472 bytes &
  \cref{tab:mesi-bill} \\
C4 & Producer--consumer model & forced writeback is 72\,B per round: 45\,\% of
  traffic with one consumer, 10\,\% with eight & \cref{eq:pc-model,fig:pc-model} \\
C5 & Superseded-writeback count & over $k$ rounds MESI writes back $k$ times,
  MOESI once & \cref{fig:wasted-timeline,eq:wb-timing} \\
C6 & One counting rule, six protocols, four patterns & ownership saves 45\,\% of
  bytes and every memory write in dirty sharing & \cref{tab:traffic-patterns,fig:bars-pc,fig:bars-mig} \\
C7 & Update versus invalidate crossover & update beats MESI below $w=10$ and
  MOESI below $w=5.5$ & \cref{fig:update-vs-invalidate}; \cite{archibald1986evaluation,thacker1987firefly} \\
C8 & Dirty remote read on point-to-point systems & DASH moves 160\,B and writes
  memory; an owner directory moves 96\,B & \cref{fig:dir-dirty-read,tab:p2p}; \cite{lenoski1990dash} \\
C9 & Five behaviours on one transaction & only the owner outcome shares dirty data
  without a memory write & \cref{fig:remote-read-tree} \\
C10 & Organisation is separate from semantics & a directory changes routing, not
  the writeback & \cref{fig:classes,tab:organisations} \\
\midrule
C11 & Controlled comparison method & MOESI as MESI plus three rules; paired runs &
  \cref{sec:method-delta,sec:method-fixed} \\
C12 & Identical cache contents & the runs differ only in writebacks &
  \cref{prop:residency}, \cref{eq:dB-writebacks} \\
C13 & Dirty Sharing Ratio & measurable with two counters under MESI & \cref{sec:dsr-def} \\
C14 & Dirty epochs and the saving identity & MOESI never moves more bytes;
  $S_B = \sigma\cdot72\,\dsr/\bar b$ & \cref{prop:never-worse}, \cref{eq:identity} \\
C15 & Effect~A and Effect~B & the latency of each, and when B hurts &
  \cref{eq:dL}, \cref{fig:two-effects} \\
C16 & Two-layer simulator & coherence checked on every access; verified against
  every analytical number & Chapter~17 \\
\bottomrule
\end{tabular}
\end{table}
